
\documentclass[11pt,reqno]{amsart}
%\usepackage[textwidth=5.85in,textheight=8.55in,centering]{geometry}
\usepackage[T1]{fontenc}
\usepackage{lmodern,amssymb,microtype}
\usepackage{xcolor}
\usepackage[colorlinks=true,linkcolor=black,citecolor=blue,urlcolor=black]{hyperref}
\hypersetup{pdftitle={A counterexample to the Erdos--Mahler conjecture},
 pdfsubject={Smooth convergents and a finite irrationality exponent},
 pdfauthor={Diego Marques}}
\newtheorem*{theorem}{Theorem}
\theoremstyle{definition}
\newtheorem*{conjecture}{Conjecture}
\newcommand{\Q}{\mathbb Q}
\newcommand{\Z}{\mathbb Z}

\theoremstyle{plain}
\newtheorem{lemma}{Lemma}[section]
\newtheorem{proposition}[lemma]{Proposition}
\newtheorem{corollary}[lemma]{Corollary}
\theoremstyle{definition}
\newtheorem{definition}[lemma]{Definition}
\theoremstyle{remark}
\newtheorem{remark}[lemma]{Remark}
\newcommand{\RR}{\mathcal R}
\newcommand{\Bad}{\mathcal A}
\newcommand{\htq}{\operatorname{ht}}
\newcommand{\norm}[1]{\lVert#1\rVert}
\newcommand{\PP}{\mathrm P}
\setlength{\emergencystretch}{1em}


\title{A counterexample to the Erd\H{o}s--Mahler conjecture}
\author[D. Marques]{Diego Marques}
\address{Departamento de Matem\'atica, Universidade de Bras\'ilia,
Bras\'ilia, DF, Brazil}
\email{diego@mat.unb.br}

\subjclass[2020]{Primary 11J70; Secondary 11J82, 11J86, 11J87}

\keywords{Erd\H{o}s--Mahler conjecture, continued fractions,
smooth convergents, irrationality exponent, Subspace Theorem,
linear forms in logarithms}

\begin{document}

\begin{abstract}
We disprove a 1939 conjecture of Erd\H{o}s and Mahler by
constructing a Cantor set of irrational numbers $\xi$ with
$4\le\mu(\xi)\le11$, each having infinitely many regular
continued-fraction convergents of the form $2^a/3^b$.
\end{abstract}

\maketitle


\section{Introduction}\label{sec:intro}


For an irrational real number $\xi$, its \emph{irrationality
exponent} $\mu(\xi)$ is the supremum of the real numbers
$\nu>0$ for which $0<|\xi-u/v|<v^{-\nu}$ holds for
infinitely many reduced fractions $u/v$. Thus $\xi$ is a
Liouville number precisely when $\mu(\xi)=\infty$.
Write $p_n/q_n$ for the regular continued-fraction
convergents of $\xi$, and $\PP(m)$ for the largest
prime factor of $|m|$, with $\PP(\pm1)=1$.

In 1939, Erd\H{o}s and Mahler \cite[p.~18]{EM}
proposed the following conjecture.

\begin{conjecture}[Erd\H{o}s--Mahler]
If $\PP(p_nq_n)$ is bounded for infinitely many $n$,
then $\xi$ is a Liouville number.
\end{conjecture}

We disprove the conjecture even when the prime factors
of the selected convergents are restricted to $2$ and $3$.

\begin{theorem}\label{thm:main}
There exists a Cantor set $K\subset\mathbb R$ such that
every $\xi\in K$ has infinitely many regular
continued-fraction convergents of the form $2^a/3^b$,
yet satisfies $4\le\mu(\xi)\le11$.
\end{theorem}

The hypothesis of the Erd\H{o}s--Mahler conjecture already
implies transcendence by Ridout's theorem \cite{Ridout};
see also \cite[Theorem~2.3, p.~29]{Bugeaud}.
Our result shows that it does not force an infinite
irrationality exponent.

Erd\H{o}s and Mahler also proved that $\xi$ is a Liouville number
if $\PP(q_{n-1}q_nq_{n+1})$ is bounded for infinitely many $n$
\cite[Theorem~2, p.~17]{EM}. Shorey~\cite{Shorey} obtained effective
refinements. The recurrence $q_{n+1}=a_{n+1}q_n+q_{n-1}$ links the
three denominators in their theorem. In the conjecture, successive
indices for which $\PP(p_nq_n)$ is bounded may be arbitrarily far apart.

Fraenkel~\cite{Fraenkel62,Fraenkel} studied rational approximation
with prescribed prime factors. Fraenkel and Borosh subsequently
obtained non-Liouville numbers by allowing the numerator and
denominator of the approximating fractions to have cofactors outside
the prescribed sets of primes. When both cofactors are bounded, their
Hausdorff dimension formula gives zero, and they explicitly left open
whether non-Liouville numbers exist in this case
\cite[p.~459]{FraenkelBorosh}. Bugeaud recorded the Erd\H{o}s--Mahler
conjecture as Problem~43 in his book~\cite[p.~215]{Bugeaud}.
More recently, Bugeaud and Nguyen~\cite[Question~4.1, p.~32]{BN}
restated it and constructed further examples of numbers with
infinitely many smooth convergents.\\

\noindent\emph{Idea of the proof.}
We choose the successive convergents from the fractions
$R(b)=2^{\lceil b\alpha\rceil}/3^b$, where $\alpha=\log 3/\log 2$.
To bound the irrationality exponent of the limit, we must control
approximation at the intervening denominators. At each stage, the
indices for which $R(b)$ lies in the prescribed interval form a set
of returns of an irrational rotation. Its threefold sumset has
cardinality comparable to that of the return set itself. Two triples
of indices with the same sum give the same product of fractions
$R(b)$. Within a suitable denominator range, sufficiently close
rational approximations inherit these identities.

For small denominators, a congruence bounds the number of possible
approximations. For large denominators, we use the inherited product
identities. At high multiplicative rank, a sumset estimate gives a
quadratic saving in the rank. At low rank, many quotients of triple
products are rational powers. A uniform height bound from the
$p$-adic Subspace Theorem allows us to choose the power exponent
separately in each denominator range. Matveev's theorem then forces
these quotients to equal the corresponding quotients of products of
the $R(b)$, and rotation discrepancy bounds their number. Summing
over the denominator ranges shows that a proportion tending to one
of the returns is admissible.

The main theorem is deduced from a counting proposition in
Section~\ref{sec:prelim}, whose proof occupies
Section~\ref{sec:rotation}. The construction is ineffective
owing to its reliance on the Subspace Theorem.

\section{The counting proposition and the main theorem}
\label{sec:prelim}

Throughout, all logarithms are natural, and
$\mathbb N=\{1,2,\ldots\}$. Rational fractions have
positive denominators. The notation $|I|$ denotes
length when $I$ is an interval and cardinality when
$I$ is a finite set. We write $F=O(G)$ or $F\ll G$
if $|F|\le CG$ for some constant $C>0$,
$F\asymp G$ if $F\ll G\ll F$, and $F=o(G)$
if $F/G\to0$. Unless otherwise stated, all
implicit constants are uniform in the varying parameters.

Set
\[
 \alpha=\frac{\log3}{\log2},\qquad
 a(b)=\lceil b\alpha\rceil,\qquad
 R(b)=\frac{2^{a(b)}}{3^b}
 \quad(b\in\mathbb N).
\]
Since $\alpha$ is irrational,
$b\alpha<a(b)<b\alpha+1$. Thus each $R(b)$ is
a reduced rational in $(1,2)$. For $t\in\mathbb N$, put
\[
 Q=3^t,\qquad A=18,\qquad N=Q^A,\qquad B=NQ^{-4},
\]
and define
\[
 \begin{gathered}
 \zeta=R(t)+Q^{-4},\qquad J=(\zeta,\zeta+Q^{-4}),\\
 \RR=\{b\in[Q,Q+N)\cap\Z:R(b)\in J\}.
 \end{gathered}
\]
Thus $|J|=Q^{-4}$ and $B=N|J|=Q^{14}$.
Since $R(t)\le2-Q^{-1}$ and $2Q^{-4}<Q^{-1}$
for $Q\ge3$, we have $J\subset(1,2)$.

A reduced rational $u/v\in[1,2]$ is
\emph{exceptional for $b$} if $u/v\ne R(b)$ and
\begin{equation}\label{eq:exceptional}
 \left|R(b)-\frac uv\right|<\frac1{16v^{11}}.
\end{equation}
Write $\RR_i$ for the returns $b\in\RR$ admitting
an exceptional fraction with $v>Q^3/6$, and put
$\RR_a=\RR\setminus\RR_i$. These returns are called
\emph{inadmissible} and \emph{admissible}, respectively.
Their dependence on $t$ is suppressed except when needed.

\begin{proposition}[Admissible returns]\label{prop:count}
As $t\to\infty$ through the positive integers,
\[
 |\RR|\asymp B,\qquad
 |\RR_i|=o(B),\qquad
 |\RR_a|\asymp B.
\]
In particular, $\RR_a$ is nonempty for every
sufficiently large $t$.
\end{proposition}

Before proving Proposition~\ref{prop:count} in
Section~\ref{sec:rotation}, we deduce the main theorem.
We first construct a single counterexample, then use
the abundance of admissible choices to obtain a Cantor set.

\begin{proof}[Proof of the theorem]
By Proposition~\ref{prop:count}, choose an integer
$t_0\ge4$ sufficiently large that $\RR_a(t)\ne\varnothing$
for every $t\ge t_0$. Recursively select
$t_{j+1}\in\RR_a(t_j)$, and put
\[
 r_j=R(t_j),\qquad Q_j=3^{t_j}\qquad(j\ge0).
\]

Since $t_{j+1}\in\RR_a(t_j)\subseteq\RR(t_j)$,
the definition of $\RR(t_j)$ gives
\begin{equation}\label{eq:successive-centers}
 Q_j^{-4}<r_{j+1}-r_j<2Q_j^{-4},
 \qquad Q_{j+1}\ge3^{Q_j}\ge2Q_j.
\end{equation}

Consider the closed intervals
\[
 K_j=[r_j+Q_j^{-4},\,r_j+3Q_j^{-4}].
\]
They lie in $(1,2)$, since $r_j\le2-Q_j^{-1}$ and
$3Q_j^{-4}<Q_j^{-1}$. Their left endpoints strictly
increase, while \eqref{eq:successive-centers} gives
\[
 r_{j+1}+3Q_{j+1}^{-4}
 <r_j+\left(2+\frac3{16}\right)Q_j^{-4}
 <r_j+3Q_j^{-4}.
\]
Hence $K_{j+1}\subset\operatorname{int}K_j$.
Moreover, $Q_j\ge2^jQ_0$, so
$|K_j|\le2Q_0^{-4}16^{-j}$. Their intersection
therefore consists of a single point $\xi\in(1,2)$,
satisfying
\begin{equation}\label{eq:limit-displacement}
 Q_j^{-4}<\xi-r_j<3Q_j^{-4}\qquad(j\ge0).
\end{equation}

If $\xi=u/v$ were rational, choose $j$ with
$Q_j^3>3v$. Since $\xi\ne r_j$, rational separation
would give
\[
 |\xi-r_j|\ge\frac1{vQ_j}>3Q_j^{-4},
\]
contrary to \eqref{eq:limit-displacement}.
Thus $\xi$ is irrational. For every $j$, $Q_j\ge3$
also gives
\[
 |\xi-r_j|<\frac3{Q_j^4}<\frac1{2Q_j^2}.
\]
Legendre's criterion
\cite[Theorem~1.8, p.~10]{Bugeaud}
therefore shows that every $r_j$ is a regular
continued-fraction convergent of $\xi$.
Their denominators are strictly increasing powers of $3$,
and their numerators are powers of $2$.
For each $0<\nu<4$, \eqref{eq:limit-displacement}
gives $0<|\xi-r_j|<Q_j^{-\nu}$ whenever
$Q_j^{4-\nu}\ge3$. Since $Q_j\ge2^jQ_0$, this holds
for infinitely many $j$, proving $\mu(\xi)\ge4$.

For the upper bound, let $u/v\in[1,2]$ be reduced,
with $v\ge Q_0$, and choose $j$ such that
\[
 Q:=Q_j\le v<Q_{j+1}=:S.
\]
Since $t_j\ge4$, we have $3^{t_j}>12t_j$, and hence
$Q^3<3^{Q/4}\le S^{1/4}$. We distinguish three ranges.

\smallskip
\noindent\emph{(i) $Q\le v\le Q^3/6$.}
If $u/v=r_j$, then $v=Q$ and
$|\xi-u/v|>v^{-4}$. Otherwise, rational separation
and \eqref{eq:limit-displacement} give
\[
 |\xi-u/v|
 \ge\frac1{vQ}-\frac3{Q^4}
 \ge\frac1{2vQ}\ge\frac1{2v^2},
\]
because $6v\le Q^3$.

\smallskip
\noindent\emph{(ii) $Q^3/6<v<S^{1/4}$.}
The reduced denominator of $r_{j+1}$ is $S>v$,
so $u/v\ne r_{j+1}$. Since
$t_{j+1}\in\RR_a(t_j)$, admissibility gives
\[
 |r_{j+1}-u/v|\ge\frac1{16v^{11}}.
\]
As $S>v^4$ and $v\ge3$,
\[
 \frac3{S^4}<\frac3{v^{16}}
 \le\frac1{32v^{11}},
\]
where the last inequality follows from
$v^5\ge3^5>96$. Applying
\eqref{eq:limit-displacement} at $j+1$, we obtain
\[
 |\xi-u/v|
 \ge\frac1{16v^{11}}-\frac3{S^4}
 \ge\frac1{32v^{11}}.
\]

\smallskip
\noindent\emph{(iii) $S^{1/4}\le v<S$.}
Again $u/v\ne r_{j+1}$. Since $S\ge3$ and $v<S$,
we have $S^3\ge9S>6v$. Rational separation yields
\[
 |\xi-u/v|
 \ge\frac1{Sv}-\frac3{S^4}
 \ge\frac1{2Sv}\ge\frac1{2v^5},
\]
using $S\le v^4$.

\smallskip
These estimates show that
\[
 |\xi-u/v|\ge\frac1{32v^{11}}
 \qquad(u/v\in[1,2],\ v\ge Q_0).
\]
For $1\le v<Q_0$, the reduced fraction $u/v$
differs from $r_0$, and $6v<Q_0^3$ gives
\[
 |\xi-u/v|
 \ge\frac1{vQ_0}-\frac3{Q_0^4}
 \ge\frac1{2vQ_0}.
\]
Also, $1+Q_0^{-1}\le r_0\le2-Q_0^{-1}$, so
\[
 K_0\subset
 \left[1+\frac1{2Q_0},\,2-\frac1{2Q_0}\right].
\]
Rationals outside $[1,2]$ consequently have distance
at least $1/(2Q_0)$ from $\xi$.
Since $Q_0\ge81$, these bounds imply
\[
 \left|\xi-\frac uv\right|
 \ge\frac1{2Q_0v^{11}}
\]
for every reduced rational $u/v$ with $v\ge1$.
In particular, $4\le\mu(\xi)\le11$.


\smallskip
\noindent\emph{A Cantor set of counterexamples.}
The preceding argument applies to every sequence of admissible
indices. For $t\in\mathbb N$, put
\[
 K(t)=[R(t)+3^{-4t},\,R(t)+3^{1-4t}].
\]
As shown above, $K(b)\subset\operatorname{int}K(t)$ whenever
$b\in\RR_a(t)$.

We claim that every sufficiently large $t$ admits two admissible
indices $b<c$ for which $K(b)$ and $K(c)$ are disjoint.
Indeed, $|\RR_a(t)|\asymp Q^{14}$, while $[Q,Q+N)$ can be
covered by $O(\log Q)$ dyadic intervals. Thus two admissible
indices satisfy $b<c<2b$. Since $R(b)$ and $R(c)$ are distinct
reduced fractions with denominators $3^b$ and $3^c$,
\[
 |R(b)-R(c)|\ge3^{-c}>3^{-2b}
 >2\cdot3^{1-4b}.
\]
As $K(b)\subset[R(b),R(b)+3^{1-4b}]$ and likewise for
$K(c)$, these intervals are disjoint.

Starting from a sufficiently large $t_0$, iterate this
selection. At level $m$ we obtain $2^m$ pairwise disjoint
closed intervals, each containing two intervals at the
next level. Their lengths are bounded by
$2Q_0^{-4}16^{-m}$, where $Q_0=3^{t_0}$.
The intersection of the resulting nested unions is
therefore a Cantor set $K$. Every $\xi\in K$ arises
from an infinite sequence of admissible indices and
hence satisfies the conclusions established above.
\end{proof}


\section{Proof of Proposition~\ref{prop:count}}
\label{sec:rotation}

The proof rests on a contrast between the size of the return
set $\RR$ and that of its threefold sumset. We shall show that
$|\RR|\asymp B$, whereas $|3\RR|\ll B$. Thus many distinct triples of return indices have
the same sum. The special
arithmetic form of $R(b)$ turns these additive relations into
exact multiplicative identities, which persist under sufficiently
close rational approximation.

Consequently, exceptional approximations in a fixed denominator
range give rise to a set of rational numbers whose threefold
product set has cardinality $O(B)$. This restriction is the
basis of the counting argument: large multiplicative rank
forces the exceptional set to be small, while small rank
will be treated through rational power classes. We first
establish the additive and multiplicative relations, then
exclude exceptional approximations in the remaining
denominator ranges.

\subsection{Many returns, few sums}

For $x\in\mathbb R$, write $\{x\}=x-\lfloor x\rfloor$ and
\[
 \norm{x}=\min_{m\in\Z}|x-m|.
\]
Thus $\norm{x}$ denotes distance to the nearest integer,
not a vector norm. For a subset $E$ of an additive group
and an integer $k\ge1$, let
\[
 kE=\{e_1+\cdots+e_k:e_i\in E\},
\]
where repetitions are allowed.

The set $\RR$ is described by an irrational rotation. Put
\[
 \Phi(y)=1-\frac{\log y}{\log2},\qquad
 I=\Phi(J),\qquad \rho=|I|.
\]
Since $\Phi(R(b))=\{b\alpha\}$, the set $\RR$ consists
precisely of the indices $b\in[Q,Q+N)\cap\Z$ for which
$\{b\alpha\}\in I$.

We use the estimate of Wu and Wang
\cite[Theorem~1]{WW}: for every $\varepsilon>0$, there
is an effective constant $\mathfrak H_\varepsilon$
such that
\begin{equation}\label{eq:wu-wang}
 |p+q_1\log2+q_2\log3|
 \ge\mathfrak H^{-4.1163051-\varepsilon},
 \qquad
 \mathfrak H=\max\{|q_1|,|q_2|\}
 \ge\mathfrak H_\varepsilon,
\end{equation}
for integers $p,q_1,q_2$.
Taking $\varepsilon=10^{-4}$, $p=0$, $q_2=-n$, and
$q_1$ nearest to $n\alpha$ in \eqref{eq:wu-wang}
gives
\[
 \log2\,\norm{n\alpha}\ge(2n)^{-4.1164051}
\]
for sufficiently large positive integers $n$.
By symmetry and the irrationality of $\alpha$,
there is $c_\kappa>0$ such that
\begin{equation}\label{eq:type}
 \norm{n\alpha}\ge c_\kappa|n|^{-\kappa}
 \quad(n\in\Z\setminus\{0\}),\qquad
 \kappa=\frac{4117}{1000},\quad
 \beta=1-\frac1\kappa=\frac{3117}{4117}.
\end{equation}

\begin{proposition}[Discrepancy]\label{lem:discrepancy}
Uniformly for integers $u$, positive integers $T$,
and circular intervals $\mathcal I\subset\mathbb R/\Z$,
\begin{equation}\label{eq:discrepancy}
 \#\{u\le n<u+T:\{n\alpha\}\in\mathcal I\}
 =T|\mathcal I|+O(T^\beta).
\end{equation}
\end{proposition}

\begin{proof}
Let $\widetilde p_i/\widetilde q_i$ denote the
convergents of $\alpha$. The standard convergent
estimate \cite[Corollary~1.4, p.~10]{Bugeaud}
and \eqref{eq:type} give
$\widetilde q_{i+1}\ll\widetilde q_i^\kappa$.

On any $\widetilde q_i$ consecutive indices, the
fractional parts of $n\alpha$ lie within
$1/\widetilde q_i$ of a translate of the uniform
grid of denominator $\widetilde q_i$. At most
$O(1)$ grid points lie sufficiently close to the
endpoints of $\mathcal I$ to affect the count.
Hence each such interval of indices has discrepancy
$O(1)$, uniformly in its initial index.

Decompose $T$ greedily into convergent-denominator
lengths. The number of blocks of length
$\widetilde q_i$ is at most
\[
 \min\{C\widetilde q_i^{\kappa-1},
          T/\widetilde q_i\}.
\]
Use the first bound for
$\widetilde q_i\le T^{1/\kappa}$ and the second
for larger denominators. Since
$\widetilde q_{i+2}\ge2\widetilde q_i$, summation
gives $O(T^{1-1/\kappa})=O(T^\beta)$.
Endpoints and intervals crossing zero contribute
only $O(1)$.
\end{proof}

The return set is large, but its threefold sumset
remains of comparable cardinality.

\begin{lemma}[Returns and their sums]\label{lem:bohr}
For all sufficiently large $t$, there is a fixed
constant $C_3>0$ such that
\begin{equation}\label{eq:return-bohr}
 |\RR|\asymp B,\qquad |3\RR|\le C_3B.
\end{equation}
\end{lemma}

\begin{proof}
Since
\[
 \frac1{2\log2}\le|\Phi'(y)|\le\frac1{\log2}
 \qquad(1\le y\le2),
\]
we have
\[
 \frac{Q^{-4}}{2\log2}\le\rho
 \le\frac{Q^{-4}}{\log2}.
\]
By \eqref{eq:discrepancy},
$|\RR|=N\rho+O(N^\beta)$.

Every $\ell\in3\RR$ lies in $[3Q,3Q+3N)$
and satisfies $\{\ell\alpha\}\in3I\pmod1$.
Since $3\rho<1$ for sufficiently large $t$,
another application of \eqref{eq:discrepancy}
gives
\[
 |3\RR|\le9N\rho+O(N^\beta).
\]
Finally,
\[
 \frac{N^\beta}{B}=Q^{4-A/\kappa}\longrightarrow0,
\]
because $A=18>4\kappa$. Since $N\rho\asymp B$,
both assertions follow.
\end{proof}

For shorter arcs, the discrepancy error may exceed
the expected number of visits. We shall also need
an estimate that retains the arc length.

\begin{lemma}[Short-arc estimate]\label{lem:short-arc}
For a circular interval $\mathcal I$ of length
$0<\lambda\le1$, any integer $u$, and any
positive integer $T$,
\begin{equation}\label{eq:short-arc}
 \#\{u\le n<u+T:\{n\alpha\}\in\mathcal I\}
 \ll T\lambda+(T^\kappa\lambda)^{1/(\kappa+1)}+1.
\end{equation}
\end{lemma}

\begin{proof}
Suppose first that $\mathcal I$ does not cross
zero. The number of visits is bounded by the
number of points of the lattice
\[
 \Lambda=
 \left\{\left(\frac nT,\frac{n\alpha-r}{\lambda}\right):
 n,r\in\Z\right\}
\]
in a translate of the closed unit square.
Its covolume is $\Delta=(T\lambda)^{-1}$.

Let $\lambda_1$ be the Euclidean length of a
shortest nonzero lattice vector. It is primitive
and extends to a lattice basis. The parallel
lattice lines are separated by $\Delta/\lambda_1$,
and their lattice points by $\lambda_1$.
Thus the square contains
\[
 O\bigl((1+\lambda_1/\Delta)
        (1+\lambda_1^{-1})\bigr)
\]
lattice points. If it contains more than one
point, then $\lambda_1\le\sqrt2$, and this
bound is $O(1+\Delta^{-1}+\lambda_1^{-1})$.

For a shortest vector with nonzero first
coordinate, \eqref{eq:type} gives
\[
 \lambda_1
 \ge\max\left\{
 \frac{|n|}{T},\frac{c_\kappa}{\lambda|n|^\kappa}
 \right\}
 \ge\left(\frac{c_\kappa}
 {T^\kappa\lambda}\right)^{1/(\kappa+1)}.
\]
For a vertical vector,
$\lambda_1\ge1/\lambda$, hence
$\lambda_1^{-1}\le1$.
Substituting these bounds proves
\eqref{eq:short-arc}. An interval crossing
zero is split into two intervals.
\end{proof}

\subsection{From additive relations to multiplicative identities}

We now use the arithmetic form of the fractions
$R(b)$. If two threefold sums of return indices
coincide, the corresponding products have the
same denominator, while their proximity forces
their numerators to coincide. The resulting
identities will also hold for sufficiently close
rational approximations.

For $X\subset\Q_{>0}$ and $k\ge1$, write
\[
 X^{[k]}=\{x_1\cdots x_k:x_i\in X\},
\]
again allowing repetitions.

\begin{lemma}[Products determined by index sums]
\label{lem:smooth-products}
For sufficiently large $t$, each $\ell\in3\RR$
determines a unique integer $h(\ell)$ such that
\begin{equation}\label{eq:smooth-product}
 Y(\ell):=\prod_{i=1}^3R(b_i)
 =\frac{2^{h(\ell)}}{3^\ell}
 \quad\text{whenever }\ell=b_1+\cdots+b_3,
 \quad b_i\in\RR.
\end{equation}
All the numbers $h(\ell)-\alpha\ell$ belong to
a common interval of length $3\rho$ contained
in $(0,3)$. Moreover, for $\ell\ne\ell'$,
\begin{equation}\label{eq:product-separation}
 |Y(\ell)-Y(\ell')|
 \ge c_\kappa\log2\,(3N)^{-\kappa}.
\end{equation}
\end{lemma}

\begin{proof}
Two products corresponding to representations
of the same $\ell$ have quotient an integral
power of $2$, and both lie in
$(\zeta^3,(\zeta+Q^{-4})^3)$. Since
\[
 \left(\frac{\zeta+Q^{-4}}{\zeta}\right)^3<2
\]
for sufficiently large $t$, their quotient
must equal one. Thus \eqref{eq:smooth-product}
defines $h(\ell)$ uniquely.

Since $a(b)-\alpha b=1-\{b\alpha\}$ and
$\{b\alpha\}\in I$, the quantities
$h(\ell)-\alpha\ell$ lie in a common
interval of length $3\rho$ contained in
$(0,3)$. For distinct $\ell,\ell'\in3\RR$,
\[
 \begin{aligned}
 |\log Y(\ell)-\log Y(\ell')|
 &=\log2\,|h(\ell)-h(\ell')
       -\alpha(\ell-\ell')|\\
 &\ge\log2\,\norm{(\ell-\ell')\alpha}\\
 &\ge c_\kappa\log2\,(3N)^{-\kappa}.
 \end{aligned}
\]
Both products exceed one, so their difference
is at least the difference of their logarithms.
This proves \eqref{eq:product-separation}.
\end{proof}

We next transfer these identities to the rational
approximations whose number must be estimated.
Fix a nonempty set $\Bad\subseteq\RR$ and a
denominator interval $(V,H]$. For each $b\in\Bad$,
choose a reduced rational $s_b\in[1,2]$, with
denominator $v_b$, satisfying
\begin{equation}\label{eq:assignment}
 V<v_b\le H,\qquad s_b\ne R(b),\qquad
 |R(b)-s_b|<(16v_b^{11})^{-1}.
\end{equation}
Put $X=\{s_b:b\in\Bad\}$ and $Z=X^{[3]}$.

\begin{lemma}[Product correspondence]\label{lem:transfer}
There exists a fixed constant $C_2>0$ such
that, for sufficiently large $t$, the conditions
\begin{equation}\label{eq:transfer-conditions}
 V^{11}>C_2N^\kappa,\qquad V^{11}>2H^6
\end{equation}
make $b\mapsto s_b$ injective and induce
a bijection
\begin{equation}\label{eq:bijection}
 \psi:3\Bad\longrightarrow Z,\qquad
 \psi(b_1+\cdots+b_3)=s_{b_1}\cdots s_{b_3}.
\end{equation}
In particular,
\begin{equation}\label{eq:product-cap-size}
 |\Bad|\le|Z|=|3\Bad|\le C_3B.
\end{equation}
If $z=U/W=\psi(\ell)$ is reduced, then
\begin{equation}\label{eq:product-accuracy}
 1\le z,Y(\ell)\le8,\qquad
 U\le8H^3,\quad W\le H^3,\qquad
 0<|Y(\ell)-z|<V^{-11}.
\end{equation}
\end{lemma}

\begin{proof}
If $s_b=s_{b'}$ for distinct $b,b'\in\Bad$,
then \eqref{eq:assignment} gives
$|R(b)-R(b')|<(8V^{11})^{-1}$.
Applying $\Phi$, we obtain
\[
 \norm{(b-b')\alpha}
 <(8\log2\,V^{11})^{-1},
\]
contrary to \eqref{eq:type}, since
$0<|b-b'|<N$ and $C_2$ is sufficiently
large. Thus $b\mapsto s_b$ is injective.

For any $b_1,\ldots,b_3\in\Bad$,
telescoping products of numbers in $[1,2]$
gives
\[
 \left|\prod_{i=1}^3R(b_i)
 -\prod_{i=1}^3s_{b_i}\right|
 \le4\sum_{i=1}^3|R(b_i)-s_{b_i}|
 <\frac34V^{-11}<V^{-11}.
\]
Equal index sums give equal smooth products
by Lemma~\ref{lem:smooth-products}.
The corresponding rational products therefore
differ by less than $2V^{-11}$.
As their reduced denominators are at most
$H^3$, rational separation and the second
condition in \eqref{eq:transfer-conditions}
force them to coincide. Hence $\psi$ is
well defined and surjective.

Conversely, if two rational products coincide,
their smooth products differ by less than
$2V^{-11}$. For distinct index sums, this
contradicts \eqref{eq:product-separation}
provided
\[
 C_2>\frac{2\cdot3^\kappa}{c_\kappa\log2}.
\]
Thus $\psi$ is injective. Fixing two summands
and applying \eqref{eq:return-bohr} proves
\eqref{eq:product-cap-size}.

The bounds on $z,U,W$ in
\eqref{eq:product-accuracy} follow by
multiplying three fractions in $[1,2]$
with denominators at most $H$.
Furthermore, rational separation gives
\[
 \frac1{3^bv_b}
 \le|R(b)-s_b|
 <\frac1{16v_b^{11}},
 \qquad
 16v_b^{10}<3^b.
\]
Consequently, if $\ell=b_1+\cdots+b_3$,
\[
 W\le\prod_{i=1}^3v_{b_i}
 <16^{-3/10}3^{\ell/10}<3^\ell.
\]
Since $Y(\ell)$ has reduced denominator
$3^\ell$, we have $z\ne Y(\ell)$.
The telescoping estimate proves the remaining
assertion of \eqref{eq:product-accuracy}.
\end{proof}

The preceding lemma bounds the size of the
threefold product set independently of the
denominator height. We shall combine this
bound with its multiplicative rank.

For a finite nonempty $X\subset\Q_{>0}$,
the \emph{affine multiplicative rank} of $X$
is the rank of the group generated by
$\{x/x_0:x\in X\}$, for any fixed $x_0\in X$.
This does not depend on the choice of $x_0$.
For a prime $p$, let $v_p(x)$ denote the
$p$-adic valuation of a positive rational $x$.

\begin{lemma}[Multiplicative rank bound]
\label{lem:rank-height}
Let $X\subset[1,2]\cap\Q$ be finite and
nonempty. Put $M=|X|$, $\mathsf P=|X^{[3]}|$,
and let $d$ be its affine multiplicative rank.
Then $X^{[3]}$ has rank $d$ and
\begin{equation}\label{eq:rank120}
 Md^2\le6\mathsf P\qquad(d\ge1).
\end{equation}
\end{lemma}

\begin{proof}
Fix $x_0\in X$ and put
\[
 E=\{(v_p(x/x_0))_p:x\in X\}.
\]
This set contains zero and spans a
$d$-dimensional real vector space.
By unique factorization, $X^{[3]}/x_0^3$
corresponds to $3E$, which has the same
rank. For $d\ge1$, the Matolcsi--Ruzsa
sumset theorem \cite[Theorem~1.5]{MR}
gives
\[
 |3E|\ge\binom{d+2}{2}
 \left(M-\frac{2d}{3}\right)
 \ge\frac{Md^2}{6},
\]
since $d\le M-1$.
\end{proof}



\subsection{Excluding small denominators}
\label{sec:lower}

We first exclude exceptional approximations with denominators
below $V_*=Q^{A/2}/t$. The position of $J$ immediately to the
right of $R(t)$ restricts the possible rational fractions to
a small number of residue classes. Combined with the short-arc
estimate, this gives a negligible contribution.

For $\nu>0$ and $0<Y_1<Y_2\le\infty$, let
$E_\nu(Y_1,Y_2)$ denote the number of indices $b\in\RR$
for which there exists a reduced rational $u/v\in[1,2]$
such that
\[
 Y_1<v\le Y_2,\qquad u/v\ne R(b),\qquad
 |R(b)-u/v|<(16v^\nu)^{-1}.
\]
Thus $|\RR_i|=E_{11}(Q^3/6,\infty)$.

\begin{lemma}[Localized fractions]
\label{lem:localized-targets}
For $V\ge Q^3/6$ and sufficiently large $Q$, at most
$10V^2/Q^4$ reduced fractions $u/v\in[1,2]$ with
$V<v\le2V$ have a radius-$(16v^9)^{-1}$ neighborhood
intersecting $J$.
\end{lemma}

\begin{proof}
Write $p=2^{a(t)}$, so that $R(t)=p/Q$ and
$\gcd(p,Q)=1$. By the definition of $J$, every such
fraction satisfies
\[
 \frac{v}{Q^3}-\frac{Q}{16v^8}
 <\Delta:=Qu-pv
 <\frac{2v}{Q^3}+\frac{Q}{16v^8}.
\]
Since $V<v\le2V$ and $V\ge Q^3/6$, these inequalities
give, for sufficiently large $Q$,
\[
 0<\Delta<\frac{5V}{Q^3}.
\]
For each integer $\Delta$ in this range, the congruence
\[
 pv\equiv-\Delta\pmod Q
\]
determines a unique residue class modulo $Q$.
It contains at most $V/Q+1\le2V/Q$ integers
$v\in(V,2V]$. Once $v$ and $\Delta$ are fixed,
so is $u=(pv+\Delta)/Q$.
The number of possible fractions is therefore at most
\[
 \frac{5V}{Q^3}\cdot\frac{2V}{Q}
 =\frac{10V^2}{Q^4}.
\]
\end{proof}

We apply this bound with approximation exponent $9$,
which also controls the exceptional set defined by
exponent $11$.

\begin{proposition}[Small denominators]
\label{prop:small-denominators}
Put $V_*=Q^{A/2}/t$. For sufficiently large $t$,
\begin{equation}\label{eq:small-denominators}
 E_{11}(Q^3/6,V_*)\ll Bt^{-2}.
\end{equation}
\end{proposition}

\begin{proof}
Fix $V\ge Q^3/6$. By
Lemma~\ref{lem:localized-targets}, there are at most
$O(Q^{-4}V^2)$ reduced fractions with denominators
in $(V,2V]$ whose radius-$(16v^9)^{-1}$
neighborhoods intersect $J$.

Under $\Phi$, each such neighborhood determines an
interval of length $O(V^{-9})$. Summing the short-arc
estimate \eqref{eq:short-arc} over these intervals gives
\[
 \begin{aligned}
 E_9(V,2V)
 &\ll NQ^{-4}V^{-7}
 +N^{\kappa/(\kappa+1)}Q^{-4}V^\sigma
 +Q^{-4}V^2,\\
 \sigma&=2-\frac9{\kappa+1}
 =\frac{1234}{5117}>0.
 \end{aligned}
\]

Sum over dyadic intervals from $Q^3/6$ to $V_*$,
truncating the last interval if necessary.
The first term decreases geometrically, while
the other two increase. The sums of the three terms are
\[
 O(BQ^{-21}),\qquad
 O(BQ^{-\vartheta A}t^{-\sigma}),\qquad
 O(Bt^{-2}),
\]
respectively, where
\[
 \vartheta=\frac{9/2-\kappa}{\kappa+1}
 =\frac{383}{5117}>0.
\]
Here we used $B=NQ^{-4}$ and
\[
 N^{-1/(\kappa+1)}V_*^\sigma
 =Q^{-\vartheta A}t^{-\sigma}.
\]
Since $A=18$ and $Q=3^t$, the first two
contributions are $O(Bt^{-2})$.
As $E_{11}\le E_9$, this proves
\eqref{eq:small-denominators}.
\end{proof}

\subsection{Exceptional heights and rational powers}
\label{sec:subspace}

It remains to control exceptional approximations with denominators
exceeding $V_*$. The first step is to show that their logarithmic
heights are $o(N)$, uniformly over the return set. This follows
from the $p$-adic Subspace Theorem. We then exploit rational
power classes among the threefold products and use Matveev's
theorem to separate pairs whose quotients do not agree with
those of the corresponding smooth products.

For $x\in\Q^\times$ and a prime $p$, write
$|x|_p=p^{-v_p(x)}$, where $v_p(-x)=v_p(x)$.
For a positive integer $n$, put
$[n]_6=2^{v_2(n)}3^{v_3(n)}$; the notation
$|\cdot|_\infty$ denotes ordinary absolute value.

We use the following two-dimensional form of the $p$-adic
Subspace Theorem. Let $\mathcal V$ be a finite set of places
of $\Q$ containing infinity, and for each
$\upsilon\in\mathcal V$ let
$L_{1,\upsilon},L_{2,\upsilon}$ be independent rational
linear forms. For fixed $c,\eta>0$, the nonzero primitive
integer vectors $\mathbf x=(x_1,x_2)$ satisfying
\begin{equation}\label{eq:subspace}
 \prod_{\upsilon\in\mathcal V}
 |L_{1,\upsilon}(\mathbf x)
 L_{2,\upsilon}(\mathbf x)|_\upsilon
 \le c\max\{|x_1|,|x_2|\}^{-\eta}
\end{equation}
lie in finitely many proper rational subspaces of $\Q^2$;
see \cite[p.~246, (1.3)--(1.5)]{Schlickewei} and
\cite[Corollary~7.2.5]{BG}. The constant $c$ is absorbed
by decreasing $\eta$ outside a bounded set.

Related Subspace-Theorem estimates for rational
approximations to quotients of power sums appear in
Corvaja and Zannier~\cite{CorvajaZannier}.
The following lemma is tailored to the family $R(b)$.

\begin{lemma}[Fixed-exponent separation]
\label{lem:fixed-separation}
For every $\varepsilon>0$, there exists
$b_0=b_0(\varepsilon)$ such that, for all integers
$b\ge b_0$ and reduced rationals $u/v\in[1,2]$
distinct from $R(b)$,
\begin{equation}\label{eq:fixed-separation}
 |R(b)-u/v|\ge3^{-\varepsilon b}v^{-2}.
\end{equation}
\end{lemma}

\begin{proof}
Write $P=2^{a(b)}$, $S=3^b$, and
$e=|Pv-Su|\ge1$. Put
\[
 \mathfrak g=\gcd(Pv,Su),\qquad
 \mathsf U=\frac{Pv}{\mathfrak g},\qquad
 \mathsf V=\frac{Su}{\mathfrak g}.
\]
The pair $(\mathsf U,\mathsf V)$ is positive and primitive.
Since $\gcd(u,v)=1$, every prime divisor of $\mathfrak g$
belongs to $\{2,3\}$. Consequently,
\[
 [\mathsf U\mathsf V]_6
 =\frac{PS[uv]_6}{\mathfrak g^2},
 \qquad
 |(\mathsf U-\mathsf V)\mathsf V|_\infty
 =\frac{eSu}{\mathfrak g^2}.
\]
It follows that
\[
 \begin{aligned}
 \Pi(\mathsf U,\mathsf V)
 &:=
 |(\mathsf U-\mathsf V)\mathsf V|_\infty
 |\mathsf U\mathsf V|_2|\mathsf U\mathsf V|_3\\
 &=\frac{eu}{P[uv]_6}
 \le2v^2|P/S-u/v|.
 \end{aligned}
\]

Suppose that \eqref{eq:fixed-separation} fails for
arbitrarily large $b$. Rational separation gives
$|P/S-u/v|\ge(Sv)^{-1}$, so any such failure
requires $v<S$. Hence, with
$\mathcal H=\max\{\mathsf U,\mathsf V\}$,
we have $\mathcal H\le2Sv<2S^2$ and
\begin{equation}\label{eq:st-inequality}
 \Pi(\mathsf U,\mathsf V)
 <2S^{-\varepsilon}
 \le2^{1+\varepsilon/2}\mathcal H^{-\varepsilon/2}.
\end{equation}

Apply \eqref{eq:subspace} to
\eqref{eq:st-inequality}, taking
$\mathcal V=\{\infty,2,3\}$, the forms
$(X-Y,Y)$ at infinity, and $(X,Y)$ at
each of the places $2$ and $3$.
The primitive pairs $(\mathsf U,\mathsf V)$
then lie on finitely many rational lines,
each containing at most one positive primitive pair.

On the other hand, along a sequence of failures
with $b\to\infty$,
\[
 0<|\mathsf U/\mathsf V-1|
 =\frac{|P/S-u/v|}{u/v}
 <S^{-\varepsilon}v^{-2}\longrightarrow0.
\]
This is impossible for a finite set of positive
primitive pairs with $\mathsf U\ne\mathsf V$.
\end{proof}

For each $t$, let $W_t$ be the maximum of $\log v$
over exceptional fractions associated with indices
$b\in\RR$, taking $W_t=0$ if there are none.
This maximum is finite: rational separation and
\eqref{eq:exceptional} give $16v^{10}<3^b$.

We claim that
\begin{equation}\label{eq:contraction}
 W_t/N\longrightarrow0\qquad(t\to\infty).
\end{equation}
Indeed, fix $\varepsilon>0$. For sufficiently large $t$,
all $b\in\RR$ satisfy $b\ge Q\ge b_0(\varepsilon)$.
Combining \eqref{eq:fixed-separation} with
\eqref{eq:exceptional} gives
$v^9<3^{\varepsilon b}/16$, and therefore
\[
 W_t\le\frac{\varepsilon\log3}{9}(Q+N).
\]
Since $Q/N\to0$, we obtain
$\limsup_{t\to\infty}W_t/N\le\varepsilon\log3/9$.
Letting $\varepsilon\downarrow0$ proves
\eqref{eq:contraction}.

We now partition the remaining denominator range.
Since $16v^{10}<3^b$ and $b<Q+N$, every exceptional
denominator is smaller than $F=3^{(Q+N)/10}$.
With $\theta=21/20$, define consecutive intervals
$(V_j,H_j]$ by
\begin{equation}\label{eq:large-blocks}
 V_0=V_*,\qquad
 H_j=\min\{V_j^\theta,F\},\qquad
 V_{j+1}=H_j,
\end{equation}
ending when $H_j=F$. There are $O(\log N)$ such
intervals, since their logarithmic endpoints
increase by the factor $\theta$ before truncation.

Let $\Bad_j$ be the set of indices $b\in\RR$
admitting an exceptional fraction with denominator
in $(V_j,H_j]$. For each such index choose the
fraction of least denominator, then least numerator.
Write $X_j$ for the selected fractions and put
\[
 Z_j=X_j^{[3]},\qquad
 \mathsf P_j=|Z_j|,
\]
with $d_j$ denoting the affine multiplicative rank
of $X_j$.

For every nonempty $\Bad_j$, the conditions of
Lemma~\ref{lem:transfer} hold uniformly:
\[
 \frac{V_j^{11}}{N^\kappa}
 \ge\frac{V_*^{11}}{N^\kappa}
 =\frac{Q^{(11/2-\kappa)A}}{t^{11}}\longrightarrow\infty,
 \qquad
 \frac{V_j^{11}}{H_j^6}
 \ge V_j^{47/10}\longrightarrow\infty.
\]
Thus the correspondence gives a bijection
$\psi_j:3\Bad_j\to Z_j$ and
$\mathsf P_j\le C_3B$.

We next count pairs of products whose quotients
are rational powers. For an integer $q\ge2$, put
\[
 \mathcal C_j(q)=
 \#\{(z,z')\in Z_j^2:
 z'/z\in(\Q_{>0}^{\times})^q\}.
\]
A pair counted by $\mathcal C_j(q)$ is
\emph{proportional} if
\[
 \frac{z'}z=\frac{Y(\ell')}{Y(\ell)},
 \qquad
 \ell=\psi_j^{-1}(z),\quad
 \ell'=\psi_j^{-1}(z').
\]
Let $\mathcal D_j(q)$ count such proportional pairs.
Both counts concern ordered pairs and include the diagonal.

\begin{lemma}[Rational power classes]\label{lem:fibers}
For every nonempty block and integer $q\ge2$,
\begin{equation}\label{eq:fibers}
 \mathcal C_j(q)=\sum_\omega n_\omega^2,\qquad
 \frac{\mathsf P_j^2}{q^{d_j}}
 \le\mathcal C_j(q)\le\mathsf P_j^2,
\end{equation}
where $n_\omega$ denotes the number of elements
of $Z_j$ in the class
$\omega\in\Q_{>0}^{\times}/(\Q_{>0}^{\times})^q$.
\end{lemma}

\begin{proof}
Fix $z_0\in Z_j$. By
Lemma~\ref{lem:rank-height}, the group generated
by the ratios $z/z_0$, $z\in Z_j$, has rank $d_j$.
Its image modulo rational $q$th powers therefore
has cardinality at most $q^{d_j}$.
The identity in \eqref{eq:fibers} counts ordered
pairs belonging to the same class. Since
$\sum_\omega n_\omega=\mathsf P_j$, the lower
bound follows from Cauchy--Schwarz; the upper
bound is immediate.
\end{proof}

The lower bound in \eqref{eq:fibers} will be
combined with an upper bound for proportional
pairs. The essential step is to show that a
nonproportional rational-power pair must have
widely separated indices. For this we need
a version of Matveev's theorem retaining the
height of the rational base in its coefficient
parameter.

For a positive rational $g=r/s$ in lowest terms,
write $\htq(g)=\log\max\{r,s\}$.

\begin{proposition}[Matveev]\label{prop:matveev}
Let $g\in\Q_{>0}\setminus\{1\}$,
$\tau=\htq(g)$, and $m,n,q\in\Z$.
If $\mathcal L=m\log2+n\log3-q\log g\ne0$, then
\begin{equation}\label{eq:matveev}
\log|\mathcal L|>
-C_0\tau(1+\log\mathfrak B),
\end{equation}
where $C_0=10^{12}$ and
$\mathfrak B=\max\{1,|q|,
\tau^{-1}\max(|m|\log2,|n|\log3)\}$.
\end{proposition}

\begin{proof}
Apply \cite[p.~1218, (1.3), and p.~1219, Corollary~2.3, (2.6)]{Matveev} over $\mathbb Q$ to the bases $2,3,g$, with height parameters $\eta_1=\log 2$, $\eta_2=\log 3$, and $\eta_3=\tau$. These parameters satisfy the requirements of \cite[(2.4)]{Matveev}, since
$\tau\ge\log 2>0.16$ and $|\log g|\le\tau$. The coefficient parameter in~(1.3) is precisely $\mathfrak B$. Since $C_1(3)\le 2^{38}$ and $2^{38}\log 2\log 3<10^{12}$, the asserted bound follows.
\end{proof}

\begin{lemma}[Separation of nonproportional pairs]
\label{lem:gap}
Put $\gamma=(4C_0)^{-1}$. For all sufficiently
large $V$, let $(V,H]=(V_j,H_j]$ be a denominator
interval, and let $q\ge2$ be an integer satisfying
\begin{equation}\label{eq:q-admissible}
 \frac{C_0(1+\log q)}q<\frac14.
\end{equation}
Every nonproportional pair counted by
$\mathcal C_j(q)$, with corresponding indices
$\ell,\ell'$, then satisfies
\begin{equation}\label{eq:gap}
 |\ell'-\ell|>(\log V)e^{\gamma q}.
\end{equation}
\end{lemma}

\begin{proof}
Write $z'/z=g^q$, $\tau=\htq(g)$,
$\delta=|\ell'-\ell|$, and $w=\log V$.
Nonproportionality and the injectivity of
$\psi_j$ imply $g\ne1$ and $\delta\ge1$.
Consider
\[
 \mathcal L=
 (h(\ell')-h(\ell))\log2
 -(\ell'-\ell)\log3-q\log g.
\]
By nonproportionality, $\mathcal L\ne0$.
Moreover, \eqref{eq:product-accuracy} and
$|(\log)'(y)|\le1$ for $y\ge1$ give
\[
 0<|\mathcal L|<2V^{-11}=2e^{-11w}.
\]
Since $|h(\ell')-h(\ell)|<2\delta+3\le5\delta$
and $\tau\ge\log2$, the bound
\eqref{eq:matveev} yields
\[
 \log|\mathcal L|>
 -C_0\tau\left(
 1+\log\max\{q,10\delta/\tau\}\right).
\]

By \eqref{eq:product-accuracy}, the reduced
numerator and denominator of $z'/z$ are at
most $8H^6$. Since $H\le V^\theta$
and $\theta=21/20$,
\[
 q\tau=\htq(z'/z)
 \le3\log2+6\log H
 \le3\log2+\frac{63}{10}w
 \le11w
\]
for sufficiently large $w$.

For fixed $q,\delta$, the function
\[
 f_\delta(y)=y\left(
 1+\log\max\{q,10\delta/y\}\right)
\]
is increasing for $y>0$: on its two
branches the derivatives are
$\log(10\delta/y)\ge\log q$ and
$1+\log q$, respectively.

Suppose that $\delta\le we^{\gamma q}$,
and put $\tau_0=11w/q$. Then
$\tau\le\tau_0$ and
\[
 \max\{q,10\delta/\tau_0\}
 \le qe^{\gamma q}.
\]
Using the monotonicity of $f_\delta$,
we obtain
\[
 \begin{aligned}
 \log|\mathcal L|
 &>-C_0f_\delta(\tau)
 \ge-C_0f_\delta(\tau_0)\\
 &\ge-11w\left(
 \frac{C_0(1+\log q)}q+C_0\gamma\right)
 >-\frac{11w}{2},
 \end{aligned}
\]
where the last inequality follows from
\eqref{eq:q-admissible} and $C_0\gamma=1/4$.
This contradicts $|\mathcal L|<2e^{-11w}$
for sufficiently large $w$, proving
\eqref{eq:gap}.
\end{proof}


\subsection{The rank dichotomy and completion}
\label{sec:adaptive}

We now choose the rational power exponent according to the
height of each denominator interval. The restriction
\eqref{eq:contraction} ensures that the separation in
Lemma~\ref{lem:gap} exceeds the entire range of possible
index sums. Thus every rational-power pair is proportional,
and the remaining count reduces to rotation discrepancy.

For each $j$ with $\Bad_j\ne\varnothing$, put
\[
 w_j=\log H_j,\qquad x_j=\log\left(\frac{N}{w_j}\right),
 \qquad K=16C_0,\quad K_1=K+1,\quad
 q_j=\lceil Kx_j\rceil.
\]
Let $\mathcal J=\{j:\Bad_j\ne\varnothing\}$ and, when
$\mathcal J\ne\varnothing$, set
$\chi=\min_{j\in\mathcal J}x_j$.
For the empty set, put $\chi=+\infty$ and interpret
all sums over $\mathcal J$ as zero.

\begin{lemma}[Proportionality at the chosen exponent]
\label{lem:proportionality}
As $t\to\infty$ through values for which
$\mathcal J\ne\varnothing$, one has $\chi\to\infty$.
For all sufficiently large $t$ and every
$j\in\mathcal J$,
\[
 \mathcal C_j(q_j)=\mathcal D_j(q_j),\qquad
 q_j\le K_1x_j,\qquad x_j<\log N.
\]
\end{lemma}

\begin{proof}
For $j\in\mathcal J$, some exceptional fraction
has denominator $v\in(V_j,H_j]$. Since
$H_j\le V_j^\theta$,
\[
 w_j\le\theta\log V_j
 <\theta\log v\le\theta W_t=o(N).
\]
By \eqref{eq:contraction},
\[
 \chi\ge\log\left(\frac{N}{\theta W_t}\right)
 \longrightarrow\infty.
\]
Thus $q_j\to\infty$ uniformly over $\mathcal J$,
and \eqref{eq:q-admissible} holds for every $q_j$
when $t$ is sufficiently large. Since
$x_j\ge\chi\ge1$, we have $q_j\le K_1x_j$;
also $H_j\ge V_*>e$ gives $x_j<\log N$.

If a pair counted by $\mathcal C_j(q_j)$ were
not proportional, Lemma~\ref{lem:gap} and
$\gamma K=4$ would give
\[
 \begin{aligned}
 |\ell'-\ell|
 &>(\log V_j)e^{\gamma q_j}
 \ge\frac{w_j}{\theta}e^{4x_j}\\
 &=\frac{N}{\theta}e^{3x_j}
 \ge\frac{N}{\theta}e^{3\chi}>3N.
 \end{aligned}
\]
But $\ell,\ell'\in3\Bad_j\subseteq[3Q,3Q+3N)$,
a contradiction. Hence every counted pair
is proportional.
\end{proof}

Proportionality imposes congruences on the exponents
of $2$ and $3$, while the rational heights bound the
difference of the corresponding indices. Together
with discrepancy, these restrictions give the following
estimate.

\begin{lemma}[Localization of proportional pairs]
\label{lem:index-bound}
For every integer $q\ge2$ and every nonempty
denominator interval $(V_j,H_j]$,
\begin{equation}\label{eq:rloc}
 \frac{\mathcal D_j(q)}{\mathsf P_j}
 \ll\frac{Q^{-4}\log(2H_j)}{q^2}
 +\left(1+\frac{\log(2H_j)}q\right)^\beta+1.
\end{equation}
\end{lemma}

\begin{proof}
Let $z=U/W$ and $z'=U'/W'$ be a proportional
pair counted by $\mathcal C_j(q)$, written
in lowest terms, with corresponding indices
$\ell,\ell'$. Then
\[
 \frac{U'W}{UW'}
 =\frac{Y(\ell')}{Y(\ell)}
 =2^{h(\ell')-h(\ell)}3^{-(\ell'-\ell)}.
\]
If $\ell'>\ell$, taking $3$-adic valuations
shows that $3^{\ell'-\ell}\mid UW'$.
Since $UW'\le8H_j^6$ by
\eqref{eq:product-accuracy}, and the
roles of the fractions can be interchanged,
\[
 |\ell'-\ell|
 \le6\log_3H_j+\log_3 8.
\]

Fix $z=\psi_j(\ell_0)$ and write
$h_0=h(\ell_0)$. Because $z'/z$ is a
rational $q$th power and
\[
 \frac{z'}z
 =2^{h(\ell')-h_0}3^{-(\ell'-\ell_0)},
\]
unique factorization gives
\[
 \ell'=\ell_0+qn,\qquad
 h(\ell')=h_0+qk
 \qquad(n,k\in\Z).
\]
The preceding height bound restricts $n$
to an interval of
$O(1+\log(2H_j)/q)$ consecutive integers.

By Lemma~\ref{lem:smooth-products}, all
$h(\ell)-\alpha\ell$ belong to one real
interval of length $3\rho$. Hence
\[
 h(\ell')-\alpha\ell'
 =h_0-\alpha\ell_0+q(k-n\alpha)
\]
restricts $k-n\alpha$ to a fixed interval
of length $3\rho/q<1$.
Each $n$ therefore determines at most one
integer $k$, while $\{n\alpha\}$ belongs
to a circular interval of length $3\rho/q$.

Applying Proposition~\ref{lem:discrepancy},
the number of possible $n$ is
\[
 O\left(
 \frac{\rho\log(2H_j)}{q^2}
 +\left(1+\frac{\log(2H_j)}q\right)^\beta
 +1\right).
\]
Each pair $(n,k)$ determines at most one
$z'$ through the bijection $\psi_j$.
Summing over the $\mathsf P_j$ choices
of $z$ and using $\rho\ll Q^{-4}$
proves \eqref{eq:rloc}.
\end{proof}

The lower bound for rational power pairs now
meets the upper bound for proportional pairs.
By Lemmas~\ref{lem:fibers} and
\ref{lem:proportionality},
\[
 \frac{\mathsf P_j^2}{q_j^{d_j}}
 \le\mathcal C_j(q_j)=\mathcal D_j(q_j).
\]
Combining this with \eqref{eq:rloc}, and using
$\log(2H_j)\le2w_j$ and
$(1+u)^\beta\le1+u^\beta$ for $u\ge0$, gives
\begin{equation}\label{eq:uniform-block}
 \mathsf P_j\ll
 Q^{-4}w_jq_j^{d_j-2}
 +w_j^\beta q_j^{d_j-\beta}
 +q_j^{d_j}.
\end{equation}

For small multiplicative rank, \eqref{eq:uniform-block}
controls $P_j$ directly. For large rank, the quadratic
saving in \eqref{eq:rank120} makes the contribution
summable over the logarithmic grid. This explains
the use of threefold products: twofold products
yield only a linear rank saving, insufficient for
this summation.

\begin{lemma}[The logarithmic grid]\label{lem:grid}
Put $s=\log\theta>0$. The numbers $x_j$,
$j\in\mathcal J$, belong to an arithmetic
grid of spacing $s$, apart from at most one
point corresponding to the final truncated
interval. If $\chi\ge1$, then for every fixed
$c>0$,
\begin{align}
 \sum_{j\in\mathcal J}e^{-cx_j}
 &\ll_{c,\theta}e^{-c\chi},
 \label{eq:grid-exp}\\
 \sum_{j\in\mathcal J}x_j^{-4/3}
 &\le2\chi^{-4/3}+\frac3s\chi^{-1/3},
 \label{eq:grid-recip}\\
 \sum_{j\in\mathcal J}e^{x_j/4}
 &\ll_\theta N^{1/4}.
 \label{eq:grid-increasing}
\end{align}
\end{lemma}

\begin{proof}
Before the final truncated interval,
$\log H_{j+1}=\theta\log H_j$, hence
$x_{j+1}=x_j-s$. The final interval
contributes at most one additional point.

The geometric series over grid points
$x_j\ge\chi$ gives \eqref{eq:grid-exp}.
For the powers $x^{-4/3}$, comparison
with an integral gives
\[
 \sum_{\substack{x\text{ on the grid}\\x\ge\chi}}
 x^{-4/3}
 \le\chi^{-4/3}
 +s^{-1}\int_\chi^\infty u^{-4/3}\,du.
\]
Adding at most $\chi^{-4/3}$ for the final
point proves \eqref{eq:grid-recip}.
Finally, $x_j<\log N$ by
Lemma~\ref{lem:proportionality}. Summing
the increasing exponential backwards from
this upper endpoint gives
\eqref{eq:grid-increasing}.
\end{proof}

We can now divide the denominator intervals
according to whether $d_j$ is smaller or
larger than $x_j^{2/3}$. In the first case,
the factor $q_j^{d_j}$ grows sufficiently
slowly for \eqref{eq:uniform-block} to be
effective. In the second, the quadratic
rank bound gives a summable estimate.

\begin{proposition}[Large denominators]
\label{prop:large-denominators}
If $\mathcal J\ne\varnothing$, then for
all sufficiently large $t$,
\begin{equation}\label{eq:large-denominators}
 \begin{aligned}
 \frac1B\sum_{j\in\mathcal J}|\Bad_j|
 \ll{}&\chi^{-1/3}+e^{-3\chi/4}\\
 &+Q^{4-A/\kappa}e^{-(\beta-1/4)\chi}
 +Q^4N^{-3/4}=o(1).
 \end{aligned}
\end{equation}
For $\mathcal J=\varnothing$, the sum is zero.
\end{proposition}

\begin{proof}
Since $K_1$ is fixed, for all sufficiently
large $x$,
\[
 x^{2/3}\log(K_1x)\le x/4.
\]
By Lemma~\ref{lem:proportionality},
$\chi\to\infty$, so this holds for all
$x_j$, $j\in\mathcal J$, when $t$ is large.

Suppose first that $d_j\le x_j^{2/3}$.
Since $q_j\le K_1x_j$,
\[
 q_j^{d_j}
 \le\exp\{x_j^{2/3}\log(K_1x_j)\}
 \le e^{x_j/4}.
\]
Using \eqref{eq:uniform-block},
$q_j\ge2$, and $|\Bad_j|\le\mathsf P_j$,
we obtain
\[
 \frac{|\Bad_j|}{B}
 \ll e^{-3x_j/4}
 +Q^{4-A/\kappa}e^{-(\beta-1/4)x_j}
 +B^{-1}e^{x_j/4},
\]
because $w_j/N=e^{-x_j}$ and
$Q^4N^{\beta-1}=Q^{4-A/\kappa}$.
Here $\beta-1/4=8351/16468>0$.
Summing over these intervals using
\eqref{eq:grid-exp} and
\eqref{eq:grid-increasing} gives
\[
 \begin{aligned}
 \frac1B
 \sum_{\substack{j\in\mathcal J\\
                 d_j\le x_j^{2/3}}}
 |\Bad_j|
 \ll{}&e^{-3\chi/4}
 +Q^{4-A/\kappa}e^{-(\beta-1/4)\chi}\\
 &+Q^4N^{-3/4},
 \end{aligned}
\]
where $B=NQ^{-4}$ was used in the last term.

If $d_j>x_j^{2/3}$, the multiplicative
rank bound \eqref{eq:rank120} and
Lemma~\ref{lem:transfer} give
\[
 |\Bad_j|
 \le\frac{6\mathsf P_j}{d_j^2}
 \le\frac{6C_3B}{d_j^2}
 <\frac{6C_3B}{x_j^{4/3}}.
\]
By \eqref{eq:grid-recip}, the sum over
these intervals is $O(B\chi^{-1/3})$.

Combining the two cases proves
\eqref{eq:large-denominators}.
Every term on its right tends to zero:
$\chi\to\infty$, $A=18>4\kappa$, and
\[
 Q^4N^{-3/4}=Q^{4-3A/4}\longrightarrow0.
\]
\end{proof}

We have now excluded all but a negligible
proportion of the returns admitting exceptional
rational approximations.

\begin{proof}[Proof of Proposition~\ref{prop:count}]
By Lemma~\ref{lem:bohr}, $|\RR|\asymp B$.
It remains to show that $|\RR_i|=o(B)$.

Every exceptional fraction associated with
$b\in\RR$ satisfies $16v^{10}<3^b$,
hence $v<F$. The intervals treated in
Proposition~\ref{prop:small-denominators}
and Proposition~\ref{prop:large-denominators}
therefore cover all inadmissible returns:
\[
 |\RR_i|=E_{11}(Q^3/6,\infty)
 \le E_{11}(Q^3/6,V_*)
 +\sum_{j\in\mathcal J}|\Bad_j|.
\]
When $\mathcal J\ne\varnothing$, the two
propositions yield
\[
 \frac{|\RR_i|}{B}
 \ll t^{-2}+\chi^{-1/3}
 +Q^{4-A/\kappa}+Q^4N^{-3/4}
 \longrightarrow0,
\]
using $e^{-3\chi/4}\ll\chi^{-1/3}$ and
$e^{-(\beta-1/4)\chi}\le1$.
When $\mathcal J=\varnothing$, the same
conclusion follows directly from
Proposition~\ref{prop:small-denominators}.

Thus $|\RR_i|=o(B)$, and consequently
\[
 |\RR_a|=|\RR|-|\RR_i|\asymp B.
\]
This completes the proof of
Proposition~\ref{prop:count}.
\end{proof}


\section*{Acknowledgments}
The author was supported by the Brazilian National Council for
Scientific and Technological Development (CNPq), Grant
No.~304467/2023-5.

\section*{AI use declaration}

The author used ChatGPT and Codex to review the exposition, check
steps in the proofs, and compare cited statements with their sources.
The author takes responsibility for the mathematical arguments and
the accuracy of the references.

\begin{thebibliography}{99}

\bibitem{BG}
E. Bombieri and W. Gubler,
\emph{Heights in Diophantine Geometry},
New Mathematical Monographs \textbf{4}, Cambridge University Press,
Cambridge, 2006.

\bibitem{Bugeaud}
Y. Bugeaud,
\emph{Approximation by Algebraic Numbers},
Cambridge Tracts in Mathematics \textbf{160}, Cambridge University Press,
Cambridge, 2004.

\bibitem{BN}
Y. Bugeaud and K. D. Nguyen,
\href{https://doi.org/10.2140/pjm.2023.326.17}{Some arithmetical properties of convergents to algebraic numbers},
\emph{Pacific J. Math.} \textbf{326} (2023), no.~1, 17--36;
\href{https://arxiv.org/abs/2209.07485}{arXiv:2209.07485}.

\bibitem{CorvajaZannier} P. Corvaja and U. Zannier,
On the length of the continued fraction for values of
quotients of power sums,
\emph{J. Th\'eor. Nombres Bordeaux} \textbf{17}
(2005), 737--748.

\bibitem{EM}
P. Erd\H{o}s and K. Mahler,
\href{https://doi.org/10.1112/jlms/s1-14.1.12}{Some arithmetical properties of the convergents of a continued fraction},
\emph{J. London Math. Soc.} \textbf{14} (1939), no.~1, 12--18.

\bibitem{Fraenkel62}
A. S. Fraenkel,
\href{https://doi.org/10.1090/S0002-9947-1962-0144852-2}{On a theorem of D. Ridout in the theory of Diophantine approximations},
\emph{Trans. Amer. Math. Soc.} \textbf{105} (1962), 84--101.

\bibitem{Fraenkel}
\bysame,
\href{https://doi.org/10.1112/jlms/s1-39.1.405}{Transcendental numbers and a conjecture of Erd\H{o}s and Mahler},
\emph{J. London Math. Soc.} \textbf{39} (1964), 405--416.

\bibitem{FraenkelBorosh} A.~S.~Fraenkel and I.~Borosh, Fractional dimension of a set of transcendental numbers, \emph{Proc. London Math. Soc.} (3) \textbf{15} (1965), 458--470. Corrigendum, \emph{ibid.} (3) \textbf{16} (1966), 192.

\bibitem{MR}
M. Matolcsi and I. Z. Ruzsa,
\href{https://doi.org/10.1007/978-0-387-68361-4_16}{Sumsets and the convex hull},
in \emph{Additive Number Theory: Festschrift in Honor of the Sixtieth Birthday of Melvyn B. Nathanson}
(D. Chudnovsky and G. Chudnovsky, eds.), Springer, New York, 2010, 221--227.

\bibitem{Matveev}
E. M. Matveev,
\href{https://doi.org/10.1070/IM2000v064n06ABEH000314}{An explicit lower bound for a homogeneous rational linear form in the logarithms of algebraic numbers. II},
\emph{Izv. Math.} \textbf{64} (2000), no.~6, 1217--1269.

\bibitem{Ridout}
D. Ridout,
\href{https://www.cambridge.org/core/journals/mathematika/article/rational-approximations-to-algebraic-numbers/9B22B10EFDE4FFC405FCAAFB12C67D3C}{Rational approximations to algebraic numbers},
\emph{Mathematika} \textbf{4} (1957), 125--131.

\bibitem{Schlickewei}
H. P. Schlickewei,
\href{https://www.numdam.org/item/CM_1992__82_3_245_0/}{The quantitative Subspace Theorem for number fields},
\emph{Compositio Math.} \textbf{82} (1992), no.~3, 245--273.

\bibitem{Shorey}
T. N. Shorey,
\href{https://doi.org/10.1016/0022-314X(83)90014-8}{Divisors of convergents of a continued fraction},
\emph{J. Number Theory} \textbf{17} (1983), no.~1, 127--133.

\bibitem{WW}
Q. Wu and L. Wang,
\href{https://doi.org/10.1016/j.jnt.2014.03.007}{On the irrationality measure of $\log3$},
\emph{J. Number Theory} \textbf{142} (2014), 264--273.
\end{thebibliography}
\end{document}
